\chapter*{附录}
\addcontentsline{toc}{chapter}{附录}
\SourcePage{179}
下面分别给出 V. Drinfeld 和 D. Gaitsgory 关于\term{nearby-cycles}{邻近轮换}、\term{vanishing-cycles}{消失轮换}函子以及 \term{fourier-transform}{Fourier 变换}函子的 \term{t-exact}{$t$-正合}性的简短证明。最初的证明分别由 Gabber（见 [Ill94]，\S\S 4.5--4.6）和 Katz--Laumon（见 [KL85]，\S 2.1.3；[Lau87]，\S 1.3.2.3）给出。

以下 $\ell$ 为素数，$\Lambda\supset\Z/\ell^n$ 为有限交换系数环，所考虑的所有\term{scheme}{概形}都定义在 $\Z[1/\ell]$ 上。

\par\medskip\noindent\textbf{A.} 设 $S$ 为一个\term{strictly-henselian}{严格 Hensel} 离散赋值环的谱，$f:X\to S$ 为\term{finite-type}{有限型} $S$-\term{scheme}{概形}。记 $s,\eta$ 为 $S$ 的闭点和泛点，$i:X_s\hookrightarrow X$、$j:X_\eta\hookrightarrow X$ 分别为闭纤维和泛纤维的嵌入。选取 $\eta$ 上可分代数的几何点 $\bar\eta/\eta$，置 $G:=\Gal(\bar\eta/\eta)$，并以 $\rho$ 表示映射 $X_{\bar\eta}\to X_\eta$。对 $\mathcal F\in\D_c^b(X_\eta,\Lambda)$，有\term{nearby-cycles}{邻近轮换}\term{complex}{复形}
\[
\psi(\mathcal F):=i^*j_*\rho_*\rho^*(\mathcal F),
\]
而对 $\mathcal G\in\D_c^b(X,\Lambda)$，有\term{vanishing-cycles}{消失轮换}\term{complex}{复形} $\phi(\mathcal G)$（见 SGA 7 XIII）。它们具有典范的\term{distinguished-triangle}{区别三角}
\[
i^*\mathcal G\longrightarrow\psi(j^*\mathcal G)\longrightarrow\phi(\mathcal G)\longrightarrow,
\]
且由 SGA $4\frac12$ 的 \emph{Théorèmes de finitude} 中定理 3.2，得 $\psi(\mathcal F),\phi(\mathcal G)\in\D_c^b(X_s,\Lambda)$。

\term{complex}{复形} $\psi(\mathcal F),\phi(\mathcal G)$ 自然具有 $G$-\term{equivariant}{等变}结构。也就是说，它们来自 $\D(X_s\times BG,\Lambda)$ 中的对象（为简便起见，仍记为 $\psi(\mathcal F),\phi(\mathcal G)$），所用函子为忘却\term{equivariant}{等变}结构函子，即沿 $X_s\to X_s\times BG$ 的拉回；这里 $X_s\times BG$ 是 $X_s$ 上带连续 $G$-作用的\term{etale}{平展}\term{sheaf}{层}所构成的\term{topos}{拓扑斯}。\footnote{事实上，SGA 7 XIII，\S 2.1 中的\term{nearby-cycles}{邻近轮换}函子为 $R\Psi_\eta:\D^+(X_\eta,\Lambda)\to\D^+(X_s\times_s\eta,\Lambda)$；在这里的\term{strictly-henselian}{严格 Hensel} 情形，$X_s\times_s\eta=X_s\times BG$。〔原注 2〕} \term{topos}{拓扑斯} $X_s\times BG$ 是 $X_s\times BG_\alpha$ 的 $2$-\term{projective}{投射}极限，其中 $G_\alpha$ 遍历 $G$ 的有限离散商。包含 $\psi(\mathcal F)$ 和 $\phi(\mathcal G)$ 的有界\term{constructible}{可构造}\term{equivariant}{等变}\term{complex}{复形}全子范畴 $\D_c^b(X_s\times BG,\Lambda)$，是各范畴 $\D_c^b(X_s\times BG_\alpha,\Lambda)$ 的 $2$-余极限。因此，对 $\D_c^b(X_s\times BG,\Lambda)$ 中的对象，其对应的 $\D_c^b(X_s,\Lambda)$ 对象上的 $G$-作用在某个开子群上平凡。

\SourcePage{180}
假设 $S$ 在某个域 $k$ 上本质\term{finite-type}{有限型}，即 $S$ 是一个\term{finite-type}{有限型} $k$-\term{scheme}{概形}的严格\term{localization}{局部化}。我们在 $\D_c^b(X,\Lambda)$ 等范畴上考虑从\term{finite-type}{有限型} $k$-\term{scheme}{概形}的情形和 $[p_{1/2}]$ 取极限而得的 \term{t-structure}{$t$-结构}；换言之，它是 [Gab04] 中对应于\term{perversity}{反常度}函数 $x\mapsto-\operatorname{tr\,deg}(k(x)/k)$ 的 \term{t-structure}{$t$-结构}在 $\D_c^b$ 上的限制。这给出各范畴 $\D_c^b(X_s\times BG,\Lambda)$ 上的 \term{t-structure}{$t$-结构}，其\term{heart}{心}为 $\D_c^b(X_s\times BG_\alpha,\Lambda)$ 上对应 \term{t-structure}{$t$-结构}的\term{heart}{心}的 $2$-余极限（这些结构可通过\term{localization}{局部化}确定，见 \S 3.2）。因此，这个\term{heart}{心}就是 $X_s$ 上带有在某个开子群上平凡的 $G$-作用的\term{perverse-sheaf}{反常层}范畴。

以下 $R\Gamma(G,-)$ 表示下列映射的全\term{direct-image}{正像}：
\[
X\times BG\longrightarrow X,\qquad X_s\times BG\longrightarrow X_s.
\]
这里的 $G$ 是 $\Z_\ell(1)$ 被一个阶与 $\ell$ 互素的 pro-群 $Q$ 的\term{extension}{扩张}。若 $\gamma$ 生成 $G/Q$，则
\[
R\Gamma(G,-):\D^+(X_s\times BG,\Lambda)\longrightarrow\D^+(X_s,\Lambda)
\]
在\term{equivariant}{等变}\term{complex}{复形}上由下式表示：
\[
K\longmapsto\operatorname{Cone}\bigl(K^Q\xrightarrow{1-\gamma}K^Q\bigr)[-1],
\]
并将 $\D_c^b$ 映入 $\D_c^b$。对 $K\in\operatorname{Perv}(X_s\times BG,\Lambda)$，置
\[
H^i(G,K):={}^{p}H^iR\Gamma(G,K)\in\operatorname{Perv}(X_s,\Lambda).
\]
于是 $H^0(G,K)=K^G$，$H^1(G,K)=K_G(-1)$；当 $i\ne0,1$ 时，$H^i(G,K)=0$。

\par\medskip\noindent\textbf{命题。}
\begin{enumerate}[label=(\roman*)]
\item 若 $\mathcal F$ 为\term{perverse-sheaf}{反常层}，则 $\psi(\mathcal F)[-1]$ 为\term{perverse-sheaf}{反常层}。
\item 若 $\mathcal G$ 为\term{perverse-sheaf}{反常层}，则 $\phi(\mathcal G)[-1]$ 为\term{perverse-sheaf}{反常层}。
\end{enumerate}
\noindent\textbf{证明。} (i) 可假设 $\psi(\mathcal F)\ne0$。设其\term{perverse-amplitude}{反常振幅}为 $[a,b]$，即
\[
{}^pH^i\psi(\mathcal F)=0\quad(i\notin[a,b]),\qquad
{}^pH^a\psi(\mathcal F),\ {}^pH^b\psi(\mathcal F)\ne0.
\]
要证 $a=b=-1$。

从 $G$ 的有限离散商所满足的相应事实取极限，得到
$\mathcal F\xrightarrow{\sim}R\Gamma(G,\rho_*\rho^*(\mathcal F))$，因而
\[
i^*j_*\mathcal F=R\Gamma\bigl(G,\psi(\mathcal F)\bigr).
\]
于是 $\psi(\mathcal F)$ 的反常\term{canonical-filtration}{典范滤过}给出一个收敛到 ${}^pH^{a+b}i^*j_*\mathcal F$ 的\term{spectral-sequence}{谱序列}，其首项为
\[
E_1^{a,b}=H^{2a+b}\bigl(G,{}^pH^{-a}\psi(\mathcal F)\bigr).
\]
由上文，$i\ne0,1$ 时 $H^i(G,-)=0$，所以该\term{spectral-sequence}{谱序列}在 $E_1$ 页退化，$i^*j_*\mathcal F$ 的\term{perverse-amplitude}{反常振幅}包含于 $[a,b+1]$，且
\begin{equation*}
\tag{$*$}
{}^pH^a i^*j_*\mathcal F=({}^pH^a\psi(\mathcal F))^G,
\qquad
{}^pH^{b+1}i^*j_*\mathcal F=({}^pH^b\psi(\mathcal F))_G(-1).
\end{equation*}
因此，若 $G$ 在 ${}^pH^a\psi(\mathcal F)$ 和 ${}^pH^b\psi(\mathcal F)$ 上的作用平凡，则 $i^*j_*\mathcal F$ 的\term{perverse-amplitude}{反常振幅}恰为 $[a,b+1]$。

\SourcePage{181}
由于 $j$ \term{affine}{仿射}，\term{complex}{复形} $j_!\mathcal F,j_*\mathcal F$ 都是\term{perverse-sheaf}{反常层}，故
\[
i^*j_*\mathcal F=i^*\operatorname{Cone}(j_!\mathcal F\to j_*\mathcal F)
\]
的\term{perverse-amplitude}{反常振幅}包含于 $[-1,0]$。据上文，若 $G$ 在 $\psi(\mathcal F)$ 上的作用平凡，结论即得证。

一般情形可通过\term{base-change}{基变换}归结为上述情形。设 $G$ 在 $\psi(\mathcal F)$ 上的作用限制到某个开子群 $G'\subset G$ 后平凡。令 $S'/S$ 为 $S$ 在 $\eta':=\bar\eta/G'$ 中的整闭包（因此，$S'$ 的剩余域是 $S$ 的剩余域的一个可能非平凡的\term{pure}{纯}不可分\term{extension}{扩张}）。将 $X/S$ 换为 $X'/S':=X\times_SS'/S'$，并将 $\mathcal F$ 换为其在 $X_{\eta'}$ 上的拉回 $\mathcal F'$。直接检查可知，$\psi(\mathcal F)$ 的拉回同构地映到 $\psi(\mathcal F')$，而后者上的 Galois 作用平凡，故命题得证。

(ii) 通过逐次约化，只需考虑 $\mathcal G$ 不可约的情形。若 $\mathcal G$ 的\term{support}{支撑}位于 $Y$ 上，结论显然；\footnote{译注：本段原文写作 $Y$；在本附录的记号中，它指闭纤维 $X_s$。} 因而可以假设
\[
\mathcal G=j_{!*}\mathcal F,\qquad\mathcal F:=j^*\mathcal G,
\]
也就是说，$\mathcal G$ 是\term{perverse-sheaf}{反常层}之间映射 $j_!\mathcal F\to j_*\mathcal F$ 的像。因此
\[
i_*i^*\mathcal G[-1]=\operatorname{Cone}(j_!\mathcal F\to\mathcal G)[-1]
\]
为\term{perverse-sheaf}{反常层}，而
\[
\operatorname{Cone}(\mathcal G\to j_*\mathcal F)
=\operatorname{Cone}(i_*i^*\mathcal G\to i_*i^*j_*\mathcal F)
\]
也为\term{perverse-sheaf}{反常层}，所以映射 $i^*\mathcal G\to i^*j_*\mathcal F$ 可以视为反常\term{truncation}{截断}
\[
({}^pH^{-1}i^*j_*\mathcal F)[1]\longrightarrow i^*j_*\mathcal F.
\]
于是由 (i) 和 $(*)$，有
\begin{align*}
\phi(\mathcal G)[-1]
&=\operatorname{Cone}\bigl({}^pH^{-1}i^*j_*\mathcal F\to{}^pH^{-1}\psi(\mathcal F)\bigr)\\
&=\operatorname{Cone}\bigl(({}^pH^{-1}\psi(\mathcal F))^G\to{}^pH^{-1}\psi(\mathcal F)\bigr).
\end{align*}
最后一个箭头是\term{perverse-sheaf}{反常层}之间显然的嵌入，故
\[
\phi(\mathcal G)[-1]={}^pH^{-1}\psi(\mathcal F)/({}^pH^{-1}\psi(\mathcal F))^G
\]
为\term{perverse-sheaf}{反常层}，证毕。\hfill$\square$

\par\medskip\noindent\textbf{注。} (i) 实际上，上述论证对任意\term{strictly-henselian}{严格 Hensel} 离散赋值环的谱 $S$ 都成立。此时须使用 $\D_c^b(X,\Lambda)$ 上与维数函数
\[
\delta_X(x)=\operatorname{tr\,deg}\bigl(k(x)/k(f(x))\bigr)+\delta_S(f(x)),
\qquad\delta_S(s)=0,\quad\delta_S(\eta)=1,
\]
相联系的反常 \term{t-structure}{$t$-结构} $[p_{1/2}]$；\footnote{这个 \term{t-structure}{$t$-结构}是 [Gab04] 中与\term{perversity}{反常度}函数 $-\delta_X$ 相联系的 \term{t-structure}{$t$-结构}在 $\D_c^b$ 上的限制。若 $S$ 是一个正规\term{finite-type}{有限型} $k$-\term{scheme}{概形}在某个不可约除子 $Y$ 的泛点处的严格\term{localization}{局部化}，则此 \term{t-structure}{$t$-结构}与前面所考虑的结构相差 $\dim Y$ 次位移。〔原注 3〕} 这时 $j_!,j_*$ 均 \term{t-exact}{$t$-正合}。更一般地，若 $X$ 是带有\term{etale}{平展}维数函数的\term{noetherian}{诺特}\term{scheme}{概形}，并且 $X$ 拟优，或在某个维数 $\le1$ 的\term{noetherian}{诺特}\term{scheme}{概形}上\term{finite-type}{有限型}，
\SourcePage{182}
则可以在 $\D_c^b(X,\Lambda)$ 上定义反常 \term{t-structure}{$t$-结构} $[p_{1/2}]$；此时，对于\term{affine}{仿射}开嵌入，$j_!,j_*$ 都 \term{t-exact}{$t$-正合}（Gabber 的定理，见 [Ill03]，\S 2.4；[PS14]，\S 1.1.2）。

(ii) 利用若干一般事实，可以从\term{torsion}{挠}系数情形推出：对系数 $\Z_\ell$，从而也对系数 $\Q_\ell$，$\psi[-1],\phi[-1]$（以及 $j_!,j_*$）关于 $[p_{1/2}]$ 和 $[p_{1/2}^+]$ 都 \term{t-exact}{$t$-正合}。在 $[p_{1/2}]$ 的情形，所用事实如下：
\begin{enumerate}[label=(\alph*)]
\item 对适当的 $X$，其上的反常 $\Z_\ell$-\term{sheaf}{层}范畴是\term{noetherian}{诺特}范畴。因此，每个\term{perverse-sheaf}{反常层}都是一个无\term{torsion}{挠}\term{perverse-sheaf}{反常层}被一个由某个 $\ell^n$ 零化的\term{perverse-sheaf}{反常层}的\term{extension}{扩张}。
\item $\D_c^b(X,\Z_\ell)$ 中的对象 $K$ 为无\term{torsion}{挠}\term{perverse-sheaf}{反常层}，当且仅当 $K\otimes\Z/\ell\in\D_c^b(X,\Z/\ell)$ 为\term{perverse-sheaf}{反常层}。
\item 函子 $\D_c^b(X,\Z/\ell^n)\to\D_c^b(X,\Z_\ell)$，即将 $\Z/\ell^n$-\term{sheaf}{层}视为 $\Z_\ell$-\term{sheaf}{层}的函子（可用 [BS15]、[Eke90] 或 [GL17] 中的任一种形式体系定义），给出范畴\term{equivalence}{等价}
\[
\operatorname{Perv}(X,\Z/\ell^n)\xrightarrow{\sim}
\operatorname{Perv}(X,\Z_\ell)_{\text{由 }\ell^n\text{ 零化}}.
\]
\end{enumerate}

(iii) 对 $G$ 的每个有限离散商 $G_\alpha$，考虑 $G_\alpha$-覆盖
\[
\rho_\alpha:X_{\eta_\alpha}\longrightarrow X_\eta,
\qquad\eta_\alpha:=\bar\eta/\Ker(G\to G_\alpha).
\]
这些 $G_\alpha$-\term{equivariant}{等变}\term{complex}{复形} $\psi_\alpha(\mathcal F)=i^*j_*\rho_{\alpha*}\rho_\alpha^*\mathcal F$ 定义了 $\D_c^b(X_s\times BG,\Lambda)$ 的一个 ind-对象。按 SGA 4 I，\S 8.4.4 的意义，这个 ind-对象本质常值，其值为 $\psi(\mathcal F)$。特别地，当 $\alpha$ 足够大时，$\psi(\mathcal F)$ 自然是 $\psi_\alpha(\mathcal F)$ 的一个直和项。

\par\medskip\noindent\textbf{B.} 假设基域 $k$ 包含有限域 $\F_q$。设 $S$ 为\term{finite-type}{有限型} $k$-\term{scheme}{概形}，$E/S$ 为秩 $n$ 的向量丛，$E^\vee$ 为其\term{duality}{对偶}，$p^\vee,p:E^\vee\times_SE\to E^\vee,E$ 为投影，$\langle\ ,\ \rangle:E^\vee\times_SE\to\mathbb A_k^1$ 为配对映射，$\psi:\F_q\to\Lambda^\times$ 为在 $\Lambda$ 的每个连通分支（或剩余域）上都非平凡的加性特征，$\mathcal L_\psi$ 为 $\mathbb A_k^1$ 上的 Artin--Schreier \term{sheaf}{层}。\term{fourier-transform}{Fourier 变换}
\[
\mathrm F_\psi:\D_c^b(E,\Lambda)\longrightarrow\D_c^b(E^\vee,\Lambda)
\]
定义为函子
\[
\mathrm F_\psi(\mathcal F):=p_!^\vee
\bigl(\langle\ ,\ \rangle^*\mathcal L_\psi\otimes p^*\mathcal F\bigr)[n].
\]

\par\medskip\noindent\textbf{命题。} $\mathrm F_\psi$ 是一个 \term{t-exact}{$t$-正合}的范畴\term{equivalence}{等价}。

\SourcePage{183}
\noindent\textbf{证明。} 我们将证明 $\mathrm F_\psi$ 是范畴\term{equivalence}{等价}，并且它与其逆函子均左 \term{t-exact}{$t$-正合}。由于 $\mathrm F_\psi$ 左\term{adjunction}{伴随}于其逆函子，由此也就得到 $\mathrm F_\psi$ 的右 \term{t-exact}{$t$-正合}性。

$\mathrm F_\psi$ 是 \term{t-exact}{$t$-正合}函子 $p^*[n]$、$\langle\ ,\ \rangle^*\mathcal L_\psi\otimes$ 与函子 $p_!$ 的复合，而 Artin 定理说明 $p_!$ 左 \term{t-exact}{$t$-正合}，故 $\mathrm F_\psi$ 左 \term{t-exact}{$t$-正合}。其余结论由下述熟知的引理推出（见 [Lau87]，\S 1.2.2.1）。

考虑 \term{fourier-transform}{Fourier 变换} $\mathrm F_{\psi^{-1}}:\D_c^b(E^\vee,\Lambda)\to\D_c^b(E,\Lambda)$。

\par\medskip\noindent\textbf{引理。} 存在自然同构
\[
\mathrm F_{\psi^{-1}}\mathrm F_\psi(\mathcal F)=\mathcal F(-n).
\]

\noindent\textbf{引理的证明。} 这里的 \term{fourier-transform}{Fourier 变换}是分别由核
\[
\langle\ ,\ \rangle^*\mathcal L_\psi[n]\in\D_c^b(E^\vee\times_SE,\Lambda),
\qquad
\langle\ ,\ \rangle^*\mathcal L_{\psi^{-1}}[n]\in\D_c^b(E\times_SE^\vee,\Lambda)
\]
定义的函子，所以复合 $\mathrm F_{\psi^{-1}}\mathrm F_\psi:\D_c^b(E,\Lambda)\to\D_c^b(E,\Lambda)$ 由核的复合 $K\in\D_c^b(E\times_SE,\Lambda)$ 定义。\footnote{若 $X,Y$ 为分离的\term{finite-type}{有限型} $S$-\term{scheme}{概形}，则每个具有有限 Tor 维数的对象 $K\in\D_c^b(Y\times_SX,\Lambda)$ 都给出函子 $\D_c^b(X,\Lambda)\to\D_c^b(Y,\Lambda)$，$\mathcal F\mapsto K(\mathcal F):=\operatorname{pr}_{Y!}(K\otimes\operatorname{pr}_X^*\mathcal F)$，称为由核 $K$ 定义的函子。若 $L\in\D_c^b(Z\times_SY,\Lambda)$ 为另一个核，则复合 $\mathcal F\mapsto L(K(\mathcal F))$ 由核的复合 $L*K:=p_{ZX!}(p_{ZY}^*L\otimes p_{YX}^*K)\in\D_c^b(Z\times_SX)$ 定义，其中 $p_{YX}:Z\times_SY\times_SX\to Y\times_SX$ 等为投影。这是分析中核的复合在\term{sheaf}{层}论中的对应物。〔原注 4〕} 我们须构造自然同构
\[
K\xrightarrow{\sim}\Delta_*(\Lambda_E)(-n),
\]
其中 $\Delta:E\hookrightarrow E\times_SE$ 为对角嵌入。$K$ 在对角线之外为零：事实上，对 $\bar k$-点 $s\in S$ 和 $(e,e')\in E_s\times E_s$，由\term{base-change}{基变换}知纤维 $K_{(e,e')}$ 等于
\[
R\Gamma_c\bigl(E_s^\vee,\langle e-e',\ \rangle^*\mathcal L_\psi\bigr)[2n].
\]
当 $e\ne e'$ 时，上式等于
\begin{align*}
&R\Gamma_c\bigl(E_s^\vee,\langle e-e',\ \rangle^!\mathcal L_\psi\bigr)(1)[2]\\
&\qquad=R\Gamma_c\bigl(\mathbb A_s^1,\mathcal L_\psi\otimes\langle e-e',\ \rangle_!\Lambda_{E^\vee}\bigr)(1-n)[2],
\end{align*}
而这等于零，因为 $\langle e-e',\ \rangle_!(\Lambda_{E^\vee})$ 为\term{constant-sheaf}{常值层}。再由\term{base-change}{基变换}，
\[
\Delta^*K=p_!(\Lambda_{E\times_SE^\vee})[2n]
=p_!p^!(\Lambda_E)(-n)=\Lambda_E(-n),
\]
证明完毕。\hfill$\square$

\par\medskip\noindent\textbf{注。} 在 \term{l-adic}{$\ell$-进}情形，即以 $\mathcal O_E$ 和 $E$ 为系数时，结论仍成立；这里 $E$ 是包含 $\mu_p$ 的 $\Q_\ell$ 有限\term{extension}{扩张}，$p=\operatorname{char}k$（可重复以上论证，或利用前面的注 (ii)）。在\term{mixed}{混合} $\overline{\Q}_\ell$ 情形，$\mathrm F_\psi$ 使\term{weight}{权}增加 $n$：事实上，若 $\mathcal F$ 是\term{weight}{权}为 $w$ 的\term{pure}{纯}\term{perverse-sheaf}{反常层}，则由 Weil II，$\mathrm F_\psi(\mathcal F)$ 的\term{weight}{权} $\le n+w$，再由引理可知其\term{weight}{权}恰为 $n+w$。
\SourcePage{184}
% 原页为空白页。
